\documentclass[11pt,a4paper]{ctexart}

\usepackage[margin=2.2cm]{geometry}
\usepackage{booktabs}
\usepackage{longtable}
\usepackage{tabularx}
\usepackage{array}
\usepackage{xcolor}
\usepackage{enumitem}
\usepackage[hidelinks]{hyperref}

\setlist[itemize]{leftmargin=2em,itemsep=2pt,topsep=3pt}
\renewcommand{\arraystretch}{1.18}
\newcolumntype{Y}{>{\raggedright\arraybackslash}X}

\title{\textbf{投资管理系统部署、测试及后续改进简要方案}}
\author{}
\date{2026年7月}

\begin{document}
\maketitle

\section{部署方案}

建议现阶段采用 Render 新加坡区域部署，先用于内部演示和小范围受控使用。项目采用 Node.js Web Service，SQLite 数据库放在 Render 持久磁盘中，域名由 Render 自动配置 HTTPS。

\subsection*{Render 费用估算}

\begin{tabularx}{\textwidth}{@{}Yrr@{}}
\toprule
项目 & 月费（美元） & 年费（美元） \\
\midrule
Hobby 工作空间 & 0 & 0 \\
Starter Web Service（512MB内存、0.5 CPU） & 7.00 & 84.00 \\
1GB 持久磁盘 & 0.25 & 3.00 \\
一个自定义域名及 HTTPS & 0 & 0 \\
\midrule
\textbf{基础合计} & \textbf{7.25} & \textbf{87.00} \\
\bottomrule
\end{tabularx}

\vspace{0.5em}
按 1 美元约合 7.2 元人民币作预算，约为每月 52 元、每年 626 元。以上不含域名、EODHD 数据、超额流量和异地备份费用。Render 免费服务不能长期保存 SQLite 数据，因此正式试用应采用付费服务和持久磁盘。价格以实际购买时为准。

\section{域名方案}

请先确认公司是否已有正式域名。若已有，建议直接使用子域名，例如
\texttt{investment.company.com} 或 \texttt{pccc.company.com}，更统一，也无需单独建立新品牌域名。

若没有公司域名，截图中首选 \texttt{puyancapital.com}，识别度和正式程度最好。域名应注册在公司或领导名下，并开启实名认证、自动续费和账号双重验证。下表的“原价”是截图中的划线参考价，不等于未来续费价；溢价域名需要特别核实续费价格。

\subsection*{截图一：puyan 系列}

\begin{longtable}{@{}lrrl@{}}
\toprule
域名 & 原价（元） & 首年价（元） & 备注 \\
\midrule
\endfirsthead
\toprule
域名 & 原价（元） & 首年价（元） & 备注 \\
\midrule
\endhead
\texttt{puyan.top}   & --  & 770.35  & 溢价域名 \\
\texttt{puyan.work}  & 68  & 14      & \\
\texttt{puyan.me}    & 108 & 38      & \\
\texttt{puyan.club}  & --  & 212.31  & 溢价域名 \\
\texttt{puyan.vip}   & --  & 1480.19 & 溢价域名 \\
\texttt{puyan.tech}  & 16  & 14      & \\
\texttt{puyan.shop}  & 238 & 10      & \\
\texttt{puyan.biz}   & --  & 212.31  & 溢价域名 \\
\texttt{puyan.cc}    & 78  & 28      & \\
\texttt{puyan.icu}   & 68  & 10      & \\
\texttt{puyan.site}  & 16  & 10      & \\
\texttt{puyan.cloud} & 178 & 14      & \\
\bottomrule
\end{longtable}

\subsection*{截图二：puyancapital 系列}

\begin{longtable}{@{}lrrl@{}}
\toprule
域名 & 原价（元） & 首年价（元） & 备注 \\
\midrule
\endfirsthead
\toprule
域名 & 原价（元） & 首年价（元） & 备注 \\
\midrule
\endhead
\texttt{puyancapital.top}  & 26  & 13 & \\
\texttt{puyancapital.com}  & 86  & 79 & 推荐；批量优惠最低价另计 \\
\texttt{puyancapital.net}  & 102 & 89 & 批量优惠最低价另计 \\
\texttt{puyancapital.xyz}  & 64  & 14 & \\
\texttt{puyancapital.work} & 68  & 14 & \\
\texttt{puyancapital.me}   & 108 & 38 & \\
\texttt{puyancapital.club} & 98  & 13 & \\
\texttt{puyancapital.ltd}  & 98  & 27 & \\
\texttt{puyancapital.vip}  & 58  & 32 & \\
\texttt{puyancapital.tech} & 16  & 14 & \\
\texttt{puyancapital.shop} & 238 & 10 & \\
\texttt{puyancapital.biz}  & 68  & 50 & \\
\texttt{puyancapital.cc}   & 78  & 28 & \\
\bottomrule
\end{longtable}

\section{人工测试分工}

\subsection*{我负责的测试}

\begin{itemize}
  \item 启动、登录、退出、角色权限和页面操作测试；
  \item 表单、文件上传、错误提示及异常输入测试；
  \item 刷新、重新登录、角色切换后的状态保存测试；
  \item 数据库保存、备份与恢复测试；
  \item CSV、PDF报告下载及页面显示测试；
  \item Chrome、Edge及不同屏幕尺寸下的基础兼容性测试；
  \item EODHD 数据更新、失败提示、更新时间及数据来源测试。
\end{itemize}

\subsection*{需要专业金融人员完成的测试}

\begin{itemize}
  \item 目标收益、CPI、无风险利率、Beta和回撤计算是否合理；
  \item 资产配置、优化限制、调仓阈值和审批规则是否符合公司制度；
  \item 收益、持仓、交易、现金流、费用、汇率及净值是否与真实记录一致；
  \item 报告指标、基准和绩效归因方法是否专业、准确且可用于正式决策；
  \item 使用一至两个历史真实案例，完整复核系统计算结果。
\end{itemize}

专业人员完成后，请将采用的规则、样例数据、预期结果和发现的问题反馈给我，我再据此调整系统。

\section{后续改进工作}

\subsection*{报告专业化}

我将负责优化报告的版式和内容，包括报告首页、核心指标、收益与风险对比、基准说明、数据日期、数据来源、图表及方法说明。绩效归因涉及配置、选股、交易执行、费用和汇率等专业口径，需要由金融专业人员先确定方法和责任划分，我再负责程序实现。

\subsection*{SQLite 正规数据库}

当前系统已经使用 SQLite，但主要以整体状态保存。后续我将把用户、会话、配置、市场数据、持仓、交易、审批、报告和审计等内容整理为独立数据库表，并补充版本迁移、备份恢复、数据约束和测试。该方案适合内部试用和小规模使用；未来若并发用户明显增加，可再迁移至 PostgreSQL。

\subsection*{EODHD 外部数据}

EODHD 不是只提供美股数据，其每日历史数据覆盖全球股票、ETF、基金、指数、外汇等多个市场。当前系统最需要每日价格、调整后价格、分红、拆股、指数/ETF、外汇和交易所信息，用于收益、Beta、回撤、回测和报告；盘中行情、期权、新闻、公司Logo等现阶段并非必需。

\begin{tabularx}{\textwidth}{@{}lrlY@{}}
\toprule
个人套餐 & 月费 & 建议 & 主要用途 \\
\midrule
Free & \$0 & 仅测试 & 每日20次、最近一年数据 \\
EOD Historical Data--All World & \$19.99 & \textbf{功能最匹配} & 全球每日历史价格、调整价格及公司行为 \\
EOD+Intraday & \$29.99 & 暂不需要 & 增加分钟和小时级盘中数据 \\
Fundamentals & \$59.99 & 暂不需要 & 财务报表、盈利及公司基本面 \\
All-In-One & \$99.99 & 暂不建议 & 行情、基本面、新闻日历及债券等综合数据 \\
\bottomrule
\end{tabularx}

\vspace{0.5em}
以上公开价格属于个人用途。若用于公司资金管理、多人内部使用、投资委员会报告或对外展示，通常应按商业用途处理，商业价格需单独询价。建议领导直接联系 EODHD，说明所需市场、代码、用户数量、每日更新频率、内部报告用途以及是否展示或导出数据，并书面确认商业授权、缓存、报告展示和数据保存权限后再购买。

\section{开发前需要确认的事项}

\begin{enumerate}[leftmargin=2em,itemsep=2pt]
  \item 公司是否已有正式域名；若没有，是否同意注册 \texttt{puyancapital.com}；
  \item 是否同意采用 Render 新加坡区域，基础预算约每年626元，不含域名和数据；
  \item 报告需要哪些正式指标和样式，以及由谁确认绩效归因方法；
  \item 由哪些金融专业人员负责业务规则和历史案例验收；
  \item EODHD 需要覆盖哪些市场、股票、ETF、指数和外汇，并由谁联系商业授权；
  \item 第一阶段是内部演示、个人受控使用，还是直接按正式多人系统建设。
\end{enumerate}

\vfill
\small
\noindent 价格参考：\href{https://render.com/pricing}{Render Pricing}、
\href{https://render.com/docs/disks}{Render Persistent Disks}、
\href{https://eodhd.com/pricing}{EODHD Pricing}、
\href{https://eodhd.com/financial-apis/commercial-vs-personal-license-use}{EODHD Commercial vs Personal Use}。价格及许可政策以购买时的官方确认结果为准。

\end{document}
